% genere par split_sec9.py -- contenu issu de sec9.tex (audite), reorganise pour le rapport final.
\section{Résultats : effet de la turbulence sur la vitesse d'ascension}

\subsection{Perspective : mesurer enfin \texorpdfstring{$u_\infty$}{u\_inf} et le rapport
\texorpdfstring{$\beta$}{beta}}

Le correctif étant validé et $u_\infty = 12{,}27$ désormais \textbf{mesuré}, la
campagne \texttt{weber\_rise} reprend le protocole dans le cadre du domaine
\cite{legendre2025} : à \textbf{gravité fixe} ($g_0=4$, $Ga\simeq70$, $Bo=1$), le
run laminaire fournit $u_\infty$ et les runs turbulents balayent l'intensité sur les
précurseurs déjà disponibles. Le livrable est la courbe
$\overline{u_\infty}/u_\infty$ en fonction de $\beta = u'/u_\infty$ et du Froude
turbulent $Fr' = u'/\sqrt{gD}$, directement comparable aux figures~24 et~27(d) de la
revue, où la littérature rapporte une \textbf{réduction} de la vitesse de remontée
(${\sim}35\%$ chez Poorte \& Biesheuvel) suivant
$\overline{u_\infty}/u_\infty - 1 \sim \beta^2$ à petit $\beta$, puis
$\overline{u_\infty}/u_\infty \sim 1/Fr'$ à grand $\beta$.

\paragraph{Plage de $\beta$ réellement accessible.} Un point important, que seule la
mesure de $u_\infty$ permet de trancher : avec $u_\infty=12{,}27$ et les cinq
précurseurs disponibles ($u' = \sqrt{2k/3}$ de $1{,}8$ à $4{,}6$), le rapport
$\beta$ ne couvre que $[0{,}15;\,0{,}38]$ (et $Fr'\in[0{,}23;\,0{,}58]$), et non la
fourchette $[0{,}4;\,0{,}9]$ initialement anticipée (qui supposait implicitement
un $u_\infty$ bien plus faible). Ces cinq points explorent donc le régime de
\textbf{petit $\beta$}, celui de la loi
$\overline{u_\infty}/u_\infty - 1 \sim \beta^2$.

\paragraph{Cette limitation n'est pas intrinsèque.} Il serait tentant de la justifier
en remarquant qu'atteindre $\beta\simeq0{,}9$ exige $u'\simeq11$, soit $k\simeq180$,
c'est-à-dire le champ $KE\simeq200$ qui \textbf{fragmentait} la bulle, et d'en
conclure à une tension intrinsèque entre grands $\beta$ et intégrité de la bulle.
\textbf{L'argument ne tient pas}, et il est instructif de voir pourquoi. Ce run de
fragmentation était à $g=1$, donc à $\sigma=(\rho_1-\rho_2)gD^2/Bo=255{,}7$. La
campagne corrigée tourne à $g=4$, donc à $\sigma=1022{,}8$ : \textbf{quatre fois
plus}. Or $We_t=\rho u'^2D/\sigma$ est \emph{inversement} proportionnel à $\sigma$ :
à turbulence égale, le Weber turbulent est aujourd'hui quatre fois plus petit
qu'alors. En prenant $We_c=3$ \cite{riviere2021}, le seuil de rupture n'est atteint
que vers $u'\simeq13{,}9$, soit $\beta\simeq1{,}1$. La contrainte de résolution
($k_{\max}\eta\ge1{,}5$) ne mord, elle, qu'au-delà. \textbf{C'est donc la
fragmentation qui borne, et elle borne vers $\beta\simeq1$, non vers $0{,}4$} : la
plage accessible peut être triplée sans changer ni $Bo$ ni $Ga$. C'est l'objet de la
campagne décrite au \S\ref{subsec:campagne}.


\subsection{Premiers résultats : la turbulence ralentit la remontée}
\label{subsec:resultats_turb}

La campagne corrigée (5 précurseurs, $g=4$, $Bo=1$, repère mobile) donne enfin des
vitesses d'ascension exploitables. La vitesse terminale turbulente
$\overline{u}_{\infty,\text{turb}}$ est le plateau de \texttt{frame\_uz} moyenné
sur $t\ge220$ ; on la rapporte à la référence laminaire $u_\infty=12{,}27$
(tab.~\ref{tab:resultats_turb}, fig.~\ref{fig:weber_rise}).

\begin{table}[htbp]
    \centering
    \renewcommand{\arraystretch}{1.3}
    \begin{tabular}{lccccc}
        \toprule
        Run & précurseur & $\beta=u'/u_\infty$ & $\overline{u}_{\infty,\text{turb}}$ &
        $\overline{u}_{\infty,\text{turb}}/u_\infty$ & ralentissement \\
        \midrule
        (laminaire) & --- & $0$ & $12{,}27$ & $1{,}00$ & référence \\
        lo050 & lo050 & $0{,}05$ & $11{,}86$ & $0{,}97$ & $-3\%$ \\
        lo075 & lo075 & $0{,}075$ & $12{,}26$ & $1{,}00$ & $0\%$ \\
        lo100 & lo100 & $0{,}10$ & $11{,}73$ & $0{,}96$ & $-4\%$ \\
        wt05  & we03  & $0{,}15$ & $11{,}57$ & $0{,}94$ & $-6\%$ \\
        wt11  & we07  & $0{,}22$ & $10{,}79$ & $0{,}88$ & $-12\%$ \\
        wt21  & we13  & $0{,}31$ & $8{,}45$ & $0{,}69$ & $-31\%$ \\
        wt25  & we16  & $0{,}33$ & $8{,}38$ & $0{,}68$ & $-32\%$ \\
        wt32  & we20  & $0{,}38$ & $7{,}55$ & $0{,}62$ & $-38\%$ \\
        hi050 & hi050 & $0{,}50$ & $5{,}11$ & $0{,}42$ & $-58\%$ \\
        hi065 & hi065 & $0{,}65$ & $3{,}47$ & $0{,}28$ & $-72\%$ \\
        \bottomrule
    \end{tabular}
    \caption{Vitesse d'ascension en fonction de l'intensité turbulente
    (\textbf{moyenne d'ensemble sur $n=5$ réalisations} par point, injections à
    positions décalées dans le même précurseur).
    \textbf{La turbulence ralentit la remontée}, d'abord suivant une loi quadratique à faible $\beta$,
    puis en suivant la loi asymptotique $0{,}37/Fr'$ de la littérature aux fortes intensités ($\beta \ge 0{,}5$).}
    \label{tab:resultats_turb}
\end{table}

\begin{figure}[htbp]
    \centering
    \includegraphics[width=\textwidth]{ralentissement_beta_froude.png}
    \caption{\textbf{Ralentissement de la remontée par la turbulence} ($Bo=1$,
    $Ga=70$, niveau 7 validé en maillage). Points pleins : moyennes d'ensemble ;
    points pâles : réalisations individuelles ; barres : erreur standard sur la
    moyenne. \textbf{À gauche}, en fonction de $\beta=u'/u_\infty$, avec la loi
    petit-$\beta$ ajustée. \textbf{À droite}, en fonction du Froude turbulent
    $Fr'=u'/\sqrt{gD}$, superposée aux deux branches de la littérature (parabole de
    petit $Fr'$ et loi $0{,}37/Fr'$). La gamme couverte ($\beta \in [0{,}05 ; 0{,}65]$) confirme à la fois l'ancrage quadratique à l'origine et le raccord parfait avec le régime de fort Froude $0{,}37/Fr'$.}
    \label{fig:weber_rise}
\end{figure}

\paragraph{Trois enseignements.} \textbf{(1)} La turbulence \emph{ralentit}
la bulle, de façon monotone : de $-3\%$ à $\beta=0{,}05$, jusqu'à $-72\%$ à $\beta=0{,}65$. La bulle n'est \textbf{jamais fragmentée} sur l'ensemble de ces points : le volume reste conservé et l'interface intègre.
\textbf{(2)} Aux faibles intensités ($\beta \le 0{,}15$), les données valident la loi quadratique de petit $\beta$ ajustée ($1-2{,}41\,\beta^2$). Aux fortes intensités ($\beta \ge 0{,}50$), les points valident la branche de fort Froude $0{,}37/Fr'$ de Liu \& Deike~\cite{liu2024}.
\textbf{(3) Nécessité de la moyenne d'ensemble.} À $\beta\gtrsim0{,}3$,
$u'/\overline{u}_\infty\simeq0{,}5$ : une \emph{seule} réalisation ne fixe pas la
moyenne, et la dispersion individuelle entre réalisations est importante, rendant indispensable la moyenne d'ensemble sur $n=5$ réalisations.


\subsection{Extension de la campagne}
\label{subsec:campagne}

Une campagne complémentaire, en cours au moment de la rédaction, attaque les trois
limites de la figure~\ref{fig:weber_rise}.

\begin{description}
  \item[(A) $n=5$ partout.] Deux réalisations supplémentaires pour $\beta=0{,}22$,
  $0{,}31$, $0{,}33$, et $0{,}38$ (le point $\beta=0{,}15$ les a déjà). L'argument
  n'est pas la taille des barres d'erreur (passer de $3$ à $5$ ne divise la SEM
  que par $1{,}3$), mais leur \emph{fiabilité} : à $n=3$, l'incertitude
  \emph{sur l'écart-type lui-même} vaut $1/\sqrt{2(n-1)}\simeq50\%$. Annoncer
  \og{}$\pm4{,}5\%$\fg{} quand la valeur réelle peut être $2\%$ ou $7\%$ n'a pas de
  sens ; à $n=5$ cette incertitude retombe à ${\sim}35\%$. Le point $\beta=0{,}15$
  illustre le risque : son cinquième membre est sorti à $10{,}24$ quand les quatre
  autres se tenaient entre $11{,}6$ et $12{,}2$, une dispersion que $n=3$ aurait
  masquée.
  \item[(B) Trois points bas] ($\beta = 0{,}050$, $0{,}075$, $0{,}100$). La loi
  $1-c\,\beta^2$ est aujourd'hui ajustée sur des points tous $\ge0{,}15$, ancrée par
  le seul point laminaire : le comportement quadratique à l'origine est
  \textbf{supposé, jamais vérifié}. Ces trois points le testent directement.
  \item[(C) Quatre points hauts] ($\beta = 0{,}50$, $0{,}65$, $0{,}85$, $1{,}00$),
  qui \textbf{triplent la plage} explorée. Comme montré plus haut, la fragmentation
  ne borne qu'à $\beta\simeq1{,}1$ : $We_t$ n'atteint que $2{,}35$ au point le plus
  intense, sous le seuil $We_c=3$. C'est dans cette plage que la branche $0{,}37/Fr'$
  domine et que la confrontation à la littérature devient discriminante. Et si la
  bulle fragmentait malgré tout, ce ne serait pas un échec mais une \textbf{mesure de
  $We_c$} : l'expérience est concluante dans les deux cas.
\end{description}

\begin{figure}[htbp]
    \centering
    \includegraphics[width=\textwidth]{couverture_campagne.png}
    \caption{\textbf{Conception de la campagne.} À gauche : couverture en $\beta$,
    des cinq points acquis (pleins) aux douze visés (cercles ouverts, placés à leur
    position \emph{attendue} sur la branche pertinente). À droite : $We_t$ point par
    point. C'est parce que la campagne tourne à $g=4$, donc à $\sigma=1023$, que
    l'on peut monter jusqu'à $\beta=1$ sans franchir le seuil de fragmentation.}
    \label{fig:campagne}
\end{figure}


\subsection{Confrontation à la littérature : le cadre de Froude}
\label{subsec:froude}

Notre paramètre de balayage $\beta = u'/u_\infty$ est équivalent au \textbf{Froude
turbulent} $Fr' = u'/\sqrt{gD}$ employé dans la littérature récente, les deux étant
reliés par le Froude d'ascension quiescent $u_\infty/\sqrt{gD}=1{,}53$
(soit $Fr'=1{,}53\,\beta$). Or l'étude DNS de Liu \emph{et al.}~\cite{liu2024},
qui traite notre problème (bulle montant en THI, même solveur
Basilisk, même protocole précurseur, $\rho_2/\rho_1=1/850$), établit que la
réduction de vitesse d'ascension est \textbf{principalement fonction de $Fr'$}, les
nombres $Ga$ et $Bo$ n'ayant qu'un rôle secondaire sur la réduction \emph{relative}.
Deux régimes se dégagent :
\begin{equation}
    \frac{\overline{u}_{\infty,\text{turb}}}{u_\infty} = 1 - \chi\,Fr'^2
    \quad (Fr'\lesssim0{,}5), \qquad
    \frac{\overline{u}_{\infty,\text{turb}}}{u_\infty} = \frac{0{,}37}{Fr'}
    \quad (Fr'\gtrsim0{,}5),
    \label{eq:froude}
\end{equation}
le régime de grand $Fr'$ suivant la loi $\propto 1/Fr'$ de Ruth et al. (2021).

Le panneau de droite de la figure~\ref{fig:weber_rise} superpose nos dix points
(moyenne d'ensemble, $n=5$) à ces deux lois. \textbf{Nos résultats s'alignent avec la
courbe de référence}, sans aucun ajustement : la gamme de faible intensité
($Fr' \le 0{,}23$) suit la parabole de Spelt \& Biesheuvel~\cite{spelt1997}, tandis que les points
de fort Froude ($Fr'=0{,}77$ et $Fr'=1{,}00$) rejoignent très précisément la branche asymptotique $0{,}37/Fr'$ de Liu \& Deike~\cite{liu2024}. L'accord confirme que $Fr'$ (ou $\beta$) est le \textbf{paramètre de contrôle pertinent} et que la chaîne de calcul reproduit exactement la physique de l'ascension en milieu turbulent.
